\section{Introduction}
\begin{figure}[t!]
        \centering
        \includegraphics[width=1\linewidth]{Figures/teaser.pdf}
         \setlength{\abovecaptionskip}{-10pt}
         \setlength{\belowcaptionskip}{-10pt}     
        \caption{\textbf{Comparison of Gaussian allocation under different target Gaussian budgets.} (a) Given the same target Gaussian budget, \OURMODEL\ allocates Gaussians non-uniformly using predicted densification scores, avoiding over-allocation in simple regions and concentrating primitives on fine-detail regions (e.g., flowers) to better preserve details even under a small budget, in contrast to voxel-based uniform allocation (AnySplat~\cite{jiang2025anysplat}). (b) Compared to pose-free and uncalibrated baselines, \OURMODEL\ achieves the best reconstruction fidelity (LPIPS/PSNR), while largely preserving the rendering quality as the number of Gaussians is reduced.}
        \label{fig:teaser}
\end{figure}
In modern computer vision,  3D scene reconstruction using deep learning has become the de facto standard. In particular, 3D Gaussian Splatting (3DGS)~\cite{kerbl20233dgs} has emerged as a highly efficient alternative to existing methodologies~\cite{mescheder2019occupancynetworks, park2019deepsdf, fridovich2022plenoxels, muller2022instant}. 3DGS represents scenes using an explicit set of 3D Gaussian primitives, enabling high-fidelity 3D reconstruction and real-time novel-view rendering.
It incorporates adaptive density control (ADC), which periodically adds or removes Gaussians during optimization. These iterative updates assign a different number of Gaussians in each region, and through this adaptive assignment, the final 3DGS representation achieves high reconstruction fidelity with a relatively small number of Gaussians.
However, the conventional 3DGS framework still inherits key limitations shared by other optimization-based 3D reconstruction methods~\cite{mildenhall2021nerf, fridovich2022plenoxels, muller2022instant}. 
It requires costly per-scene iterative optimization and typically relies on densely captured input views with known camera parameters, which can be impractical in real-world scenarios. 
This has motivated feed-forward 3DGS methods~\cite{charatan2024pixelsplat, chen2024mvsplat, xu2025depthsplat, ye2024nopo, li2025vicasplat, huang2025spfsplat, jiang2025anysplat}, which are trained on large-scale datasets to build strong 3D priors. These frameworks can reconstruct a 3D scene from only a few input images through a single forward pass, preserving the real-time rendering capabilities of 3DGS while enabling generalization to unseen scenes. 
However, existing feed-forward 3DGS methods have a significant limitation in that they do not allocate Gaussians efficiently. This limitation stems from removing the iterative optimization process of conventional 3DGS, which also eliminates its periodic adaptive density control (ADC) that densifies Gaussians during training.
Most works~\cite{charatan2024pixelsplat, chen2024mvsplat, xu2025depthsplat, ye2024nopo, li2025vicasplat, zhang2025flare, huang2025spfsplat} adopt a pixel-to-Gaussian pipeline, which assigns Gaussians at the pixel level. This fixes the total number of Gaussians to the number of pixels in the entire image and prevents flexible adjustment of Gaussian positions, resulting in duplicated Gaussians across different views. AnySplat~\cite{jiang2025anysplat}, which employs a voxel-to-Gaussian pipeline, can adjust the number of Gaussians by changing the voxel size, but this typically requires training a new model. Moreover, because it allocates Gaussians uniformly in space (i.e., assigning one Gaussian per voxel), it struggles to produce a high-quality and compact Gaussian representation under a limited Gaussian budget.

To address inefficient Gaussian allocation, we introduce \textbf{\OURMODEL}, a feed-forward network that performs \emph{predictive densification} for 3D Gaussian Splatting from a set of uncalibrated images.
Our approach predicts densification decisions in a single forward pass, treating Gaussian densification as a learnable prediction problem within a unified feed-forward pipeline.
Specifically, the network estimates a \emph{densification score} that indicates whether additional Gaussians should be allocated to each region.
By estimating both spatial complexity and multi-view overlap, the predicted densification score avoids over-allocation in simple regions and prevents duplicate Gaussians in areas covered by overlapping input images.
This feed-forward densification strategy enables spatially adaptive allocation and yields compact Gaussian representations while maintaining competitive reconstruction fidelity. As illustrated in~\cref{fig:teaser}, this allows \OURMODEL\ to concentrate Gaussians on fine-detail regions while avoiding unnecessary allocation in simple regions, achieving higher rendering quality under the same Gaussian budget.
Through extensive experiments, \OURMODEL\ achieves on-par or superior novel-view synthesis quality while using significantly fewer Gaussians than prior uncalibrated feed-forward methods that rely solely on image inputs.

The contributions of our work can be summarized as:
\begin{itemize}
    \item \textbf{Gaussian-count controllable feed-forward 3DGS.} We propose \OURMODEL, a feed-forward framework that reconstructs 3D Gaussian Splatting representations from sparse, uncalibrated images while enabling explicit control over the final number of Gaussians through feed-forward predictive densification.

    \item \textbf{Densification-score-guided allocation for high fidelity under a limited budget.} We introduce a densification score that predicts where additional Gaussians should be allocated, enabling spatially adaptive Gaussian allocation without iterative optimization and maintaining high reconstruction fidelity even under a limited Gaussian budget.

    \item \textbf{State-of-the-art performance in the uncalibrated setting.} Extensive experiments show that \OURMODEL\ achieves on-par or superior novel-view synthesis quality while using significantly fewer Gaussians than prior uncalibrated feed-forward methods that rely solely on image inputs.
\end{itemize}